\section*{Capítulo \ref{ch:g6:circuits-and-safety} --- Circuitos eléctricos y seguridad}

\begin{solution}{exo:g6:circuits-and-safety:1}
Con fidelidad: las conexiones --- qué aparato se une con cuál y dónde
se bifurcan las carreteras. Ignora: la longitud, el color y el
serpenteo real de los cables sobre la mesa.
\end{solution}

\begin{solution}{exo:g6:circuits-and-safety:2}
\omterm{def:g3:first-electric-circuit:battery}{Pila}: dos barras paralelas desiguales; lámpara: un círculo con una
cruz; \omterm{ex:g3:first-electric-circuit:switch}{interruptor} abierto: un puente levadizo levantado;
\omterm{ex:g3:first-electric-circuit:switch}{interruptor} cerrado: el puente caído; \omterm{ex:g6:circuits-and-safety:motor}{motor}: una M dentro de
un círculo; \omterm{ex:g6:circuits-and-safety:motor}{zumbador}: un círculo atravesado. La barra \emph{larga}
marca el $+$ de la \omterm{def:g3:first-electric-circuit:battery}{pila}.
\end{solution}

\begin{solution}{exo:g6:circuits-and-safety:3}
Un rectángulo: \omterm{def:g3:first-electric-circuit:battery}{pila}, pulsador y \omterm{ex:g6:circuits-and-safety:motor}{zumbador} en un único
bucle. Al levantar el dedo se abre el \omterm{ex:g3:first-electric-circuit:switch}{interruptor} --- un hueco en
la única carretera ---, y la ley del bucle calla el \omterm{ex:g6:circuits-and-safety:motor}{zumbador} al
instante.
\end{solution}

\begin{solution}{exo:g6:circuits-and-safety:4}
No necesariamente. La pregunta que decide: ¿quién está conectado con
quién? Los mismos aparatos con las mismas conexiones --- el mismo
circuito, por muy distintas que parezcan las marañas.
\end{solution}

\begin{solution}{exo:g6:circuits-and-safety:5}
El ventilador sigue girando: la muerte de la lámpara apaga solo su
propia \omterm{def:g5:series-parallel-circuits:parallel}{rama} --- la independencia del paralelo. El
\omterm{ex:g3:first-electric-circuit:switch}{interruptor} abierto está en la carretera común, antes de la
bifurcación, así que para a los dos: la regla del puesto de mando.
\end{solution}

\begin{solution}{exo:g6:circuits-and-safety:6}
Una carretera conductora desnuda que une directamente los dos lados de
la fuente, saltándose todos los aparatos: la corriente corre sin
estorbos y los cables se calientan. Víctimas: la \omterm{def:g3:first-electric-circuit:battery}{pila}
(agotada y arruinada), el aislamiento del cable --- y, en la red de la
casa, a veces la casa entera.
\end{solution}

\begin{solution}{exo:g6:circuits-and-safety:7}
Ha cortado su \omterm{def:g5:series-parallel-circuits:parallel}{rama} en un parpadeo porque la corriente empezó a
desbocarse --- una avería, no un capricho. Hay que encontrar y arreglar
la avería (desenchufar al sospechoso, revisar la \omterm{def:g5:series-parallel-circuits:parallel}{rama}) antes de
volver a subir al guardia.
\end{solution}

\begin{solution}{exo:g6:circuits-and-safety:8}
Bañera: la piel húmeda conduce, y el empuje de la red a través de un
cuerpo mojado puede ser mortal. Cables dañados: un aislamiento
agrietado es una valla con un agujero --- los dedos pueden encontrarse
con el cobre. \omterm{def:g3:first-electric-circuit:battery}{Pila} sí, enchufe jamás: el empuje de la \omterm{def:g3:first-electric-circuit:battery}{pila} es
diminuto; el del enchufe es cientos de veces más tremendo.
\end{solution}

\begin{solution}{exo:g6:circuits-and-safety:9}
Esquema: \omterm{def:g3:first-electric-circuit:battery}{pila}, \omterm{ex:g3:first-electric-circuit:switch}{interruptor} cerrado en la carretera común y después
una bifurcación --- \omterm{def:g5:series-parallel-circuits:parallel}{rama} del \omterm{ex:g6:circuits-and-safety:motor}{zumbador} y \omterm{def:g5:series-parallel-circuits:parallel}{rama} de la
lámpara --- que se vuelven a juntar hacia el $-$. Predicción: la
lámpara alumbra \emph{y} el \omterm{ex:g6:circuits-and-safety:motor}{zumbador} suena, de manera independiente.
Con un \omterm{ex:g3:first-electric-circuit:switch}{interruptor} extra abierto en la \omterm{def:g5:series-parallel-circuits:parallel}{rama} del \omterm{ex:g6:circuits-and-safety:motor}{zumbador}:
silencio ahí, pero la lámpara sigue alumbrando --- los
\omterm{ex:g3:first-electric-circuit:switch}{interruptores} de \omterm{def:g5:series-parallel-circuits:parallel}{rama} mandan solo sobre su \omterm{def:g5:series-parallel-circuits:parallel}{rama}.
\end{solution}

\begin{solution}{exo:g6:circuits-and-safety:10}
El cable extra es un \omterm{def:g6:circuits-and-safety:short}{cortocircuito}: una carretera desnuda y directa
del $+$ al $-$. La corriente prefiere la carretera fácil y vacía al
\omterm{ex:g6:circuits-and-safety:motor}{motor}, que trabaja duro, así que el ventilador pasa hambre
mientras el cable atajo carga con la avalancha y se calienta ---
agotando la \omterm{def:g3:first-electric-circuit:battery}{pila} e invitando a las quemaduras. Fuera el cable de la
``firmeza''.
\end{solution}

\begin{solution}{exo:g6:circuits-and-safety:11}
El \omterm{ex:g3:first-electric-circuit:switch}{interruptor} de pared abre el bucle propio de la lámpara --- la
carretera de la lámpara queda cortada. El magnetotérmico abre la
\omterm{def:g5:series-parallel-circuits:parallel}{rama} entera más arriba --- aunque la instalación de la lámpara
esté averiada o el \omterm{ex:g3:first-electric-circuit:switch}{interruptor} de pared esté mal etiquetado, la
\omterm{def:g5:series-parallel-circuits:parallel}{rama} está muerta. Dos huecos \omterm{def:g4:series-circuits:series}{en serie}: cualquiera de los dos
por sí solo para la corriente, y los dos juntos perdonan una
equivocación sobre cualquiera de ellos.
\end{solution}

\begin{solution}{exo:g6:circuits-and-safety:12}
\omterm{def:g3:first-electric-circuit:battery}{Pila}, después el \omterm{ex:g3:first-electric-circuit:switch}{interruptor} general en la carretera común y
después cuatro \omterm{def:g5:series-parallel-circuits:parallel}{ramas} paralelas: foco 1; foco 2; \omterm{ex:g6:circuits-and-safety:motor}{motor} del
telón con su propio \omterm{ex:g3:first-electric-circuit:switch}{interruptor} \omterm{def:g4:series-circuits:series}{en serie}; \omterm{ex:g6:circuits-and-safety:motor}{zumbador} con su
propio \omterm{ex:g3:first-electric-circuit:switch}{interruptor} \omterm{def:g4:series-circuits:series}{en serie}. Comprobaciones: el general manda sobre
todo (carretera común); los focos se funden de manera independiente
(\omterm{def:g5:series-parallel-circuits:parallel}{ramas} separadas); el \omterm{ex:g6:circuits-and-safety:motor}{motor} y el \omterm{ex:g6:circuits-and-safety:motor}{zumbador} obedecen cada
uno a su \omterm{ex:g3:first-electric-circuit:switch}{interruptor} (\omterm{def:g4:series-circuits:series}{en serie} dentro de su \omterm{def:g5:series-parallel-circuits:parallel}{rama}) y se ignoran
mutuamente (\omterm{def:g5:series-parallel-circuits:parallel}{en paralelo} entre \omterm{def:g5:series-parallel-circuits:parallel}{ramas}).
\end{solution}

\begin{solution}{pb:g6:circuits-and-safety:1}
\textbf{1.} \omterm{def:g4:series-circuits:series}{En serie}, un aparato fundido o apagado dejaría todo a
oscuras, y tres aparatos se repartirían el empuje, todos flojos y
enclenques. El paralelo mantiene cada aparato independiente y a plena
potencia.
\textbf{2.} Cada \omterm{ex:g3:first-electric-circuit:switch}{interruptor} manda exactamente sobre el aparato de
su \omterm{def:g5:series-parallel-circuits:parallel}{rama}; ninguno manda sobre las otras \omterm{def:g5:series-parallel-circuits:parallel}{ramas}.
\textbf{3.} En la carretera común, entre la \omterm{def:g3:first-electric-circuit:battery}{pila} y la primera
bifurcación --- aguas arriba de todas las \omterm{def:g5:series-parallel-circuits:parallel}{ramas}.
\textbf{4.} \omterm{def:g3:first-electric-circuit:battery}{Pila} $\to$ \omterm{ex:g3:first-electric-circuit:switch}{interruptor} general $\to$ \omterm{def:g5:series-parallel-circuits:parallel}{nudo} que
se bifurca en tres \omterm{def:g5:series-parallel-circuits:parallel}{ramas} --- (lámpara de techo $+$ su
\omterm{ex:g3:first-electric-circuit:switch}{interruptor}), (lámpara de banco $+$ su \omterm{ex:g3:first-electric-circuit:switch}{interruptor}),
(\omterm{ex:g6:circuits-and-safety:motor}{motor} del ventilador $+$ su \omterm{ex:g3:first-electric-circuit:switch}{interruptor}) --- que se vuelven
a juntar y regresan a la \omterm{def:g3:first-electric-circuit:battery}{pila}.
\textbf{5.} Están \omterm{def:g4:series-circuits:series}{en serie} en una misma \omterm{def:g5:series-parallel-circuits:parallel}{rama}: se reparten el
empuje (los dos enclenques --- la lámpara de techo flojita era en
realidad el síntoma de la pareja del banco), y el \omterm{ex:g3:first-electric-circuit:switch}{interruptor}
del ventilador, que está en su carretera compartida, manda sobre los
dos.
\textbf{6.} Separarlos: tender un cable nuevo para que la lámpara del
banco tenga su propia \omterm{def:g5:series-parallel-circuits:parallel}{rama} sobre la alimentación --- así cada uno
queda a plena potencia y es independiente.
\textbf{7.} Un \omterm{def:g6:circuits-and-safety:short}{cortocircuito} entre los \omterm{def:g3:first-electric-circuit:battery}{bornes}. Hay que
comprobar si la \omterm{def:g3:first-electric-circuit:battery}{pila} sigue dando corriente --- los \omterm{def:g6:circuits-and-safety:short}{cortocircuitos}
la agotan y pueden arruinarla (una \omterm{def:g3:first-electric-circuit:battery}{pila} caliente merece la
jubilación).
\textbf{8.} Los \omterm{def:g6:circuits-and-safety:short}{cortocircuitos} (cables desnudos que se encuentran
entre sí) y las descargas o quemaduras (dedos que se encuentran con el
cobre desnudo).
\textbf{9.} La piel húmeda conduce muchísimo mejor que la seca; hasta
el montaje con \omterm{def:g3:first-electric-circuit:battery}{pilas} merece la disciplina de las manos secas --- y
en la casa ese mismo contacto se encuentra con el empuje de la red,
cientos de veces más tremendo: posiblemente mortal. Las costumbres hay
que entrenarlas en el circuito seguro.
\textbf{10.} El magnetotérmico ha detectado que por su \omterm{def:g5:series-parallel-circuits:parallel}{rama}
corría demasiada corriente --- una avería, probablemente lluvia dentro
de una lámpara. Orden sensato: dejar el magnetotérmico bajado,
desenchufar y revisar la \omterm{def:g5:series-parallel-circuits:parallel}{rama}, arreglar o secar la avería y solo
entonces volver a subirlo.
\textbf{11.} Su propio bucle, separado de la red de la caseta (o como
una \omterm{def:g5:series-parallel-circuits:parallel}{rama} paralela más sobre la \omterm{def:g3:first-electric-circuit:battery}{pila} de la caseta):
\omterm{def:g3:first-electric-circuit:battery}{pila}, pulsador en la casa, cable doble largo,
\omterm{ex:g6:circuits-and-safety:motor}{zumbador} en la caseta y vuelta --- un timbre es un bucle \omterm{def:g4:series-circuits:series}{en serie}
propio por muy lejos que se estiren sus cables.
\textbf{12.} Por ejemplo: ``Enchufes: nuestros experimentos acaban en
la \omterm{def:g3:first-electric-circuit:battery}{pila} --- el enchufe no es nuestro. Agua: las manos mojadas no
tocan ninguna instalación, ni dentro ni fuera. Cables desnudos: todas
las uniones forradas; un cable desnudo es una avería, no un atajo''.
\end{solution}
